\chapter{Biết Nhịn Mình Đúng Lúc}
\markboth{Biết Nhịn Mình Đúng Lúc}{Knowing When to Restrain Yourself}

\section{Chiếc Xe Đẩy Ở Siêu Thị}
\begin{truyen}{Giữa Cần Và Muốn Chỉ Cách Nhau Một Bàn Tay}{Need and Want Are Separated Only by a Hand}
\chuhoa{C}{ó những buổi đi siêu thị không ồn ào nhưng đủ làm lộ hết thói quen tiêu tiền của cả nhà. Mỗi quầy hàng đều có một lý do nghe rất hợp lý để bỏ thêm một món vào xe đẩy.}
\textit{(There are shopping trips that are not dramatic, yet they reveal a family’s spending habits with surprising clarity. Every shelf offers a reason that sounds reasonable enough to add one more item to the cart.)}

Con muốn đồ chơi, người lớn muốn đồ tiện, còn chính người cầm tiền cũng có vài thứ thích mắt. Nếu không tỉnh, người ta rất dễ gọi tất cả là cần thiết. Nhưng khi về đến quầy thanh toán, những cái ``cũng nên mua'' ấy bỗng hợp lại thành một khoản nặng vai.
\textit{(Children want toys, adults want convenience, and even the one holding the money finds a few things appealing. Without alertness, everything begins to look necessary. But at the cashier, all those ``might as well buy it'' decisions gather into a heavy total.)}

Người thầy dừng lại, chia giấy thành hai cột: cần và muốn. Chỉ một thao tác nhỏ ấy thôi đã cứu được cả buổi mua sắm. Bởi khi gọi đúng tên một món, sức hút của nó yếu đi rất nhiều.
\textit{(The teacher stopped and divided a sheet of paper into two columns: need and want. That small act alone rescued the entire shopping trip. For when something is named accurately, much of its power fades.)}
\end{truyen}

\begin{baihoc}
Nhiều khi tiết kiệm không cần tài năng lớn; nó chỉ cần sự gọi tên cho đúng.
\textit{(Often saving requires no great talent; it simply requires naming things correctly.)}
\end{baihoc}

\begin{ghinhoanh}
Name it clearly.

Needs and wants must not wear the same mask.
\end{ghinhoanh}

\begin{tuhocnhanh}
1. Em đang gọi món ``muốn'' nào là ``cần''? \textit{Which desire are you calling a need?}

2. Nếu chia hai cột hôm nay, em sẽ bỏ lại món gì? \textit{If you made two columns today, what would you leave behind?}
\end{tuhocnhanh}
\ngancach

\section{Kéo Mạ Cho Mau Lớn}
\begin{truyen}{Câu Chuyện Của Sự Sốt Ruột}{The Story of Impatience}
\chuhoa{M}{ạnh Tử kể chuyện người nông dân thấy mạ chậm lớn nên nhổ lên cho cao thêm. Ông ta tưởng mình đang giúp, nhưng cuối cùng cả ruộng mạ héo đi.}
\textit{(Mencius told of a farmer who thought his seedlings were growing too slowly, so he pulled them upward to make them taller. He believed he was helping, but the whole field withered.)}

Trong chuyện tiền bạc, nhiều người cũng làm điều tương tự với đời mình. Họ muốn được thoải mái ngay, muốn sắm đủ ngay, muốn giải tỏa mệt mỏi ngay. Cái gì chậm một chút là thấy khó chịu. Nhưng đời sống bền không lớn lên bằng cách chiều mình liên tục.
\textit{(Many people do something similar with their financial lives. They want comfort immediately, possessions immediately, relief from tiredness immediately. Anything slow becomes irritating. But a stable life does not grow by constant self-indulgence.)}

Biết nhịn mình đúng lúc không phải là tàn nhẫn với bản thân. Đó là hiểu rằng có những điều phải chờ thêm một chút mới giữ được lâu. Sự sốt ruột trong tiêu tiền thường trả giá đắt hơn ta nghĩ.
\textit{(Knowing when to restrain oneself is not cruelty toward the self. It is understanding that some things must wait a little longer if they are to last. Impatience in spending usually costs more than we think.)}
\end{truyen}

\begin{baihoc}
Không phải thứ gì đến chậm cũng xấu. Nhiều điều đến chậm để ta đủ sức giữ nó lâu hơn.
\textit{(Not everything that comes slowly is bad. Many things come slowly so that we may be strong enough to keep them longer.)}
\end{baihoc}

\begin{ghinhoanh}
Impatience is expensive.

Some good things must arrive slowly.
\end{ghinhoanh}

\begin{tuhocnhanh}
1. Điều gì em đang muốn có ngay nhưng thật ra có thể chờ? \textit{What do you want immediately that can actually wait?}

2. Sự sốt ruột nào đang làm em tiêu tiền sai? \textit{What impatience is causing you to spend wrongly?}
\end{tuhocnhanh}
\ngancach

\section{Một Lần Trả Món Về Kệ}
\begin{truyen}{Niềm Vui Ngắn Bị Đổi Lấy Sự Nhẹ Lòng Dài Hơn}{A Short Pleasure Exchanged for a Longer Relief}
\chuhoa{K}{hi trả vài món về lại kệ, cả nhà không ai vui ngay. Nhưng tối hôm đó, nhìn số tiền còn giữ được, người thầy thấy lòng mình nhẹ hơn.}
\textit{(When a few items were returned to the shelf, no one felt pleased immediately. But that evening, seeing how much money had been preserved, the teacher felt lighter inside.)}

Đó là một bài học rất đời thường: có những niềm vui phải bỏ qua trong chốc lát để đổi lấy sự yên tâm dài hơn. Người không quen cảm giác bị từ chối ngay trong hiện tại thường sẽ phải gặp cảm giác lo âu dữ hơn ở tương lai.
\textit{(It was a very practical lesson: some pleasures must be passed over for a moment in order to gain a longer peace of mind. Those who cannot tolerate small denials in the present often face greater anxiety in the future.)}

Tiết kiệm vì thế không phải lúc nào cũng vui. Nhưng nó có một phần thưởng lặng hơn và sâu hơn: cảm giác mình không buông tay với cuộc đời của chính mình.
\textit{(Saving is therefore not always pleasant. But it offers a quieter and deeper reward: the feeling that one has not let go of one’s own life.)}
\end{truyen}

\begin{baihoc}
Một lần nhịn đúng có thể đổi lấy nhiều ngày nhẹ lòng sau đó.
\textit{(One wise act of restraint can buy many days of relief afterward.)}
\end{baihoc}

\begin{ghinhoanh}
Short denial, longer peace.

That is one face of wisdom.
\end{ghinhoanh}

\begin{tuhocnhanh}
1. Em có dám bỏ qua một niềm vui ngắn để giữ sự an tâm dài hơn không? \textit{Can you give up a short pleasure for a longer peace?}

2. Lần gần nhất em tự thắng được mình là khi nào? \textit{When was the last time you overcame yourself?}
\end{tuhocnhanh}

\begin{turanminh}
1. Tôi phải tập gọi đúng tên mọi ham muốn: cái gì là cần thì mua, cái gì chỉ là muốn thì phải biết chờ.

2. Không nhịn được mình ở siêu thị, ngoài quán hay trên điện thoại thì rất khó đòi hỏi một tương lai tài chính yên ổn.

3. Một lần bớt mua đúng lúc có thể không vui ngay, nhưng nó cứu được nhiều ngày sau khỏi nặng lòng.
\end{turanminh}

\begin{dayhoc}
1. Đây là chương rất phù hợp để dạy học sinh phân biệt ``cần'' và ``muốn'' bằng những ví dụ gần với đời sống của các em.

2. Thầy có thể yêu cầu các em liệt kê ba món mình muốn mua nhất, rồi tự xếp lại xem món nào có thể chờ một tuần, một tháng, hay lâu hơn.

3. Bài học nên nhấn mạnh là trưởng thành không phải muốn gì cũng có ngay, mà là biết chờ điều đáng có.
\end{dayhoc}